\documentclass[aspectratio=169]{beamer}
\usetheme{Madrid}
\usecolortheme{seahorse}

\usepackage{graphicx}
\usepackage{booktabs}
\usepackage{amsmath}
\usepackage{xcolor}
\usepackage{tabularx}
\usepackage{array}

\title[Chapter 1: Picturing Distributions]{\textbf{Chapter 1: Picturing Distributions with Graphs}}
\subtitle{School of Mathematical and Computational Sciences \\ University of Prince Edward Island}
\author{Ali Raisolsadat}
\date{Part I: Exploring Data}

\setbeamertemplate{navigation symbols}{}
\setbeamertemplate{itemize items}[circle]
\setbeamertemplate{enumerate items}[default]

% Looks for figures in the current directory or in a figures/ subfolder.
\graphicspath{{./}{figures/}}

% This helper lets the slides compile even before the RMarkdown file has
% generated its PNG figures. Once the files exist, the real figures appear.
\newcommand{\lectureimage}[2][0.98\linewidth]{%
  \IfFileExists{#2}{%
    \includegraphics[width=#1]{#2}%
  }{%
    \fbox{\parbox[c][0.29\textheight][c]{0.92\linewidth}{%
      \centering\small
      Figure file: \texttt{\detokenize{#2}}\\[0.15cm]
      Generate this figure by knitting the revised Chapter 1 RMarkdown notes.
    }}%
  }%
}

\begin{document}

% =============================================================
% PART I: INTRODUCTION TO STATISTICS AND WHY GRAPHS MATTER
% =============================================================
\section{Why Statistics?}

% SLIDE 1
\begin{frame}
  \titlepage
\end{frame}

% SLIDE 2
\begin{frame}{The Course Roadmap: The 5-Part Journey}
  \begin{columns}[T]
    \begin{column}{0.48\textwidth}
      \textbf{I: Exploring Data}
      \begin{itemize}
        \item Ch 1--2: Distributions \& Numerical Summaries
        \item Ch 3: Scatterplots \& Correlation
        \item Ch 4: Regression
      \end{itemize}
      \vspace{0.1cm}
      \textbf{II: Probability \& Distributions}
      \begin{itemize}
        \item Ch 9--10: Probability Rules
        \item Ch 11: Normal Distributions
        \item Ch 12: Binomial / Poisson
      \end{itemize}
      \vspace{0.1cm}
      \textbf{III: Data Generation}
      \begin{itemize}
        \item Ch 6: Samples \& Observational Studies
        \item Ch 7: Designing Experiments
      \end{itemize}
    \end{column}
    \begin{column}{0.48\textwidth}
      \textbf{IV: Foundations of Inference}
      \begin{itemize}
        \item Ch 13: Sampling Distributions
        \item Ch 14: Confidence Intervals \& $p$-values
        \item Ch 15: Inference in Practice
      \end{itemize}
      \vspace{0.1cm}
      \textbf{V: Specific Inference Procedures}
      \begin{itemize}
        \item Ch 17: Inference about a Population Mean
        \item Ch 19: Inference about a Population Proportion
      \end{itemize}
      \vspace{0.2cm}
      \begin{block}{Where We Are Today}
        Chapter 1 begins with a simple question: \textbf{What does the data look like?}
      \end{block}
    \end{column}
  \end{columns}
\end{frame}

% SLIDE 3
%\begin{frame}{Chapter 1 Learning Objectives}
%  By the end of this chapter, you should be able to:
%  \begin{enumerate}
%    \item identify \textbf{individuals} and \textbf{variables} in a dataset;
%    \item distinguish \textbf{categorical} and \textbf{quantitative} variables;
%    \item recognize \textbf{discrete} and \textbf{continuous} quantitative variables;
%    \item choose an appropriate \textbf{bar graph, pie chart, histogram, dotplot, or time plot};
%    \item describe a quantitative distribution using \textbf{SOCS};
%    \item recognize common ways graphs can be \textbf{misleading}; and
%    \item use graphs to detect possible \textbf{outliers and data-quality problems}.
%  \end{enumerate}
%\end{frame}

% SLIDE 4
\begin{frame}{What Is Statistics?}
  \begin{block}{Definition}
    \textbf{Statistics is the science of learning from data.}
  \end{block}
  \begin{block}{Ernest Rutherford}
    \textbf{All science is either physics or stamp collecting.}
  \end{block}
  \vspace{0.15cm}
  Statistics gives us tools to:
  \begin{itemize}
    \item collect information;
    \item organize and display data;
    \item summarize patterns, compare groups and use data to answer questions.
  \end{itemize}
  \vspace{0.15cm}
  \begin{alertblock}{Statistics Is More Than Formulas}
    A statistical thinker asks: \emph{What does the data show, what does it not show and how confident should we be in the conclusion?}
  \end{alertblock}
\end{frame}

% SLIDE 5
\begin{frame}{Statistics Is Everywhere}
  \begin{itemize}
    \item \textbf{Biology:} Do plants given fertilizer grow taller than plants that do not receive it?
    \item \textbf{Chemistry:} Does changing the catalyst increase reaction yield?
    \item \textbf{Psychology:} Do students who sleep longer tend to report better concentration?
    \item \textbf{Economics:} Are grocery prices increasing over time, or was it only COVID that distort our impression?
    \item \textbf{Artificial Intelligence:} Does the training dataset fairly represent the population where the AI system will be used?
  \end{itemize}
\end{frame}

% SLIDE 6
\begin{frame}{Why Statistics Still Matters}
  \vspace{0.15cm}
  \begin{columns}[T]
    \begin{column}{0.47\textwidth}
      \begin{block}{LLMs and AI Can Help Us}
        \begin{itemize}
          \item organize data quickly;
          \item calculate summaries;
          \item generate visualizations;
          \item identify possible patterns.
        \end{itemize}
      \end{block}
    \end{column}
    \begin{column}{0.49\textwidth}
      \begin{alertblock}{But We Still Must Ask}
        \begin{itemize}
          \item Is this the right graph?
          \item Is the scale fair?
          \item Are there errors or unusual values?
          \item Does the conclusion actually follow from the data?
        \end{itemize}
      \end{alertblock}
    \end{column}
  \end{columns}
\end{frame}

% SLIDE 7
\begin{frame}{A First Lesson: Averages Do Not Tell the Whole Story}
  \begin{columns}[c]
    \begin{column}{0.56\textwidth}
      \centering
      \lectureimage{fig_same_mean_different_data.png}
    \end{column}
    \begin{column}{0.41\textwidth}
      Both samples have an average plant height of \textbf{20 cm}.
      \begin{itemize}
        \item \textbf{Sample A:} tightly clustered near 20 cm.
        \item \textbf{Sample B:} much more variable.
      \end{itemize}
      \begin{block}{Key Idea}
        A single number can hide important structure. \textbf{Graphs help us see what summaries miss.}
      \end{block}
    \end{column}
  \end{columns}
\end{frame}

% =============================================================
% PART II: INDIVIDUALS AND VARIABLES
% =============================================================
\section{Individuals and Variables}

% SLIDE 8
\begin{frame}{Anatomy of a Dataset}
  \begin{block}{The Rectangular Data Idea}
    In a typical dataset \textbf{rows represent individuals or cases} and \textbf{columns represent variables}.
  \end{block}
  \vspace{0.15cm}
  \begin{table}
    \centering
    \small
    \begin{tabular}{lcccc}
      \toprule
      \textbf{Plant ID} & \textbf{Fertilizer} & \textbf{Germinated?} & \textbf{Leaves} & \textbf{Height (cm)} \\
      \midrule
      P01 & None & Yes & 5 & 12.4 \\
      P02 & None & Yes & 6 & 13.1 \\
      P03 & Fertilizer A & Yes & 8 & 16.7 \\
      P04 & Fertilizer A & No & 0 & 0.0 \\
      P05 & Fertilizer B & Yes & 7 & 15.9 \\
      \bottomrule
    \end{tabular}
  \end{table}
  \begin{alertblock}{First Question for Any Dataset}
    \textbf{What does one row represent?}
  \end{alertblock}
\end{frame}

% SLIDE 9
\begin{frame}{Individuals, Variables and Identifiers}
  From the plant dataset:
  \begin{itemize}
    \item \textbf{Individuals:} the individual plants.
    \item \textbf{Variables:}
    \begin{itemize}
      \item Fertilizer Group
      \item Germinated?
      \item Number of Leaves
      \item Height
    \end{itemize}
    \item \textbf{Plant ID:} mainly an identifier used to keep track of each plant.
  \end{itemize}
  \vspace{0.2cm}
  \begin{block}{Definition}
    A \textbf{variable} is a characteristic recorded for each individual that can take different values.
  \end{block}
\end{frame}

% SLIDE 10
\begin{frame}{Two Fundamental Types of Variables}
  \begin{columns}[T]
    \begin{column}{0.48\textwidth}
      \begin{block}{Qualitative (Categorical)}
        Places an individual into a group or category.
        \begin{itemize}
          \item Blood type: A, B, AB, O
          \item Fertilizer group: None, A, B
          \item Reaction outcome: Complete, Partial, None
          \item University program
        \end{itemize}
      \end{block}
    \end{column}
    \begin{column}{0.48\textwidth}
      \begin{block}{Quantitative (Numerical)}
        Records a numerical amount for which arithmetic is meaningful.
        \begin{itemize}
          \item Plant height (cm)
          \item Reaction time (ms)
          \item Temperature ($^\circ$C)
          \item Mass (g)
          \item Price (\$)
        \end{itemize}
      \end{block}
    \end{column}
  \end{columns}
\end{frame}

% SLIDE 11
\begin{frame}{Categorical Variables: Nominal vs. Ordinal}
  \begin{columns}[T]
    \begin{column}{0.48\textwidth}
      \begin{block}{Nominal}
        Categories have \textbf{no natural order}.
        \begin{itemize}
          \item Blood type
          \item Treatment group
          \item Reaction outcome type
          \item Academic program
        \end{itemize}
      \end{block}
    \end{column}
    \begin{column}{0.48\textwidth}
      \begin{block}{Ordinal}
        Categories have a \textbf{meaningful order}.
        \begin{itemize}
          \item Low / Medium / High
          \item Poor / Fair / Good / Excellent
          \item Mild / Moderate / Severe
        \end{itemize}
      \end{block}
    \end{column}
  \end{columns}
  \vspace{0.2cm}
  \centering
  \textbf{The labels describe groups; the labels themselves are not measurements.}
\end{frame}

% SLIDE 12
\begin{frame}{The Coding Trap: Numbers Can Still Be Categories}
  Suppose a spreadsheet uses:
  \[
    1 = \text{Base Fertilizer}, \qquad 2 = \text{Fertilizer A}, \qquad 3 = \text{Fertilizer B}.
  \]
  \begin{itemize}
    \item The values are written using numbers.
    \item But the numbers are only \textbf{codes for categories}.
    \item An average treatment code of $1.8$ has no useful scientific meaning.
  \end{itemize}
  \begin{alertblock}{Diagnostic Question}
    \textbf{Would taking an average of these numbers have a meaningful interpretation?}
    \begin{itemize}
      \item Yes $\rightarrow$ usually quantitative.
      \item No $\rightarrow$ usually categorical.
    \end{itemize}
  \end{alertblock}
\end{frame}

% SLIDE 13
\begin{frame}{Quantitative Variables: Discrete vs. Continuous}
  \begin{columns}[T]
    \begin{column}{0.48\textwidth}
      \begin{block}{Discrete: Usually Counting}
        \begin{itemize}
          \item Number of leaves
          \item Number of bacterial colonies
          \item Number of customers
          \item Number of incorrect answers
        \end{itemize}
      \end{block}
    \end{column}
    \begin{column}{0.48\textwidth}
      \begin{block}{Continuous: Usually Measuring}
        \begin{itemize}
          \item Plant height
          \item Body mass
          \item Reaction time
          \item Temperature
          \item Concentration
        \end{itemize}
      \end{block}
    \end{column}
  \end{columns}
\end{frame}

% SLIDE 14
\begin{frame}{Count or Measure?}
  \begin{table}
    \centering
    \begin{tabular}{lll}
      \toprule
      \textbf{Question} & \textbf{Usually Produces} & \textbf{Example} \\
      \midrule
      How many? & Discrete quantitative & Number of leaves \\
      How much? & Continuous quantitative & Plant mass \\
      Which group? & Categorical & Fertilizer type \\
      What level/rating? & Ordinal categorical & Low / Medium / High \\
      \bottomrule
    \end{tabular}
  \end{table}
  \vspace{0.25cm}
  \begin{block}{Good Starting Point}
    This rule is not perfect for every advanced application but it works well for most introductory examples.
  \end{block}
\end{frame}

% SLIDE 15
\begin{frame}{From Variable Type to Graph Type}
  \begin{table}
    \centering
    \small
    \begin{tabularx}{\textwidth}{>{\raggedright\arraybackslash}X >{\raggedright\arraybackslash}X >{\raggedright\arraybackslash}X}
      \toprule
      \textbf{Situation} & \textbf{Good Starting Graph} & \textbf{Main Question} \\
      \midrule
      One categorical variable & Bar graph & Which categories are most or least common? \\
      Parts of one whole & Pie chart & How is one whole divided? \\
      Small quantitative sample & Dotplot & What are the individual values and overall pattern? \\
      Larger quantitative sample & Histogram & What is the shape of the distribution? \\
      Measurements over time & Time plot & Is there a trend, cycle, or sudden change? \\
      \bottomrule
    \end{tabularx}
  \end{table}
  \vspace{0.1cm}
  \centering
  \textbf{The best graph depends on both the variable type and the question.}
\end{frame}

% =============================================================
% PART III: CATEGORICAL VARIABLES
% =============================================================
\section{Categorical Data}

% SLIDE 16
\begin{frame}{The Distribution of a Categorical Variable}
  The \textbf{distribution} tells us what values a variable takes and how often they occur.
  \vspace{0.15cm}
  For categorical data, we summarize with:
  \begin{itemize}
    \item \textbf{Frequency (count):} number of observations in each category.
    \item \textbf{Relative frequency (proportion):}
    \[
      \text{Proportion} = \frac{\text{Count}}{n}.
    \]
    \item \textbf{Percentage:} proportion $\times 100\%$.
  \end{itemize}
  \begin{block}{Example}
    If 18 of 40 students choose biology as their preferred science elective,
    \[
      \frac{18}{40}=0.45=45\%.
    \]
  \end{block}
\end{frame}

% SLIDE 17
\begin{frame}{Example 1.1: Reaction Outcomes in a Chemistry Lab}
  Suppose 50 reaction trials are classified as follows:
  \begin{table}
    \centering
    \begin{tabular}{lrr}
      \toprule
      \textbf{Reaction Outcome} & \textbf{Count} & \textbf{Percentage} \\
      \midrule
      Complete reaction & 28 & 56\% \\
      Partial reaction & 15 & 30\% \\
      No visible reaction & 7 & 14\% \\
      \midrule
      \textbf{Total} & \textbf{50} & \textbf{100\%} \\
      \bottomrule
    \end{tabular}
  \end{table}
  \begin{block}{Table or Graph?}
    A table gives exact values. A graph makes the overall pattern easier to see quickly.
  \end{block}
\end{frame}

% SLIDE 18
\begin{frame}{Bar Graphs: The Default Choice for Categorical Data}
  A \textbf{bar graph} represents each category with a separate bar.
  \begin{itemize}
    \item Bar height can show \textbf{count, proportion, or percentage}.
    \item Bars are \textbf{separated} because the categories are distinct.
    \item Nominal categories can often be reordered to make comparisons easier.
    \item Ordinal categories should usually keep their natural order.
  \end{itemize}
  \vspace{0.15cm}
  \begin{alertblock}{When Reading a Bar Graph}
    Which category is largest? Smallest? Are any similar? Are the values counts or percentages? Does the vertical axis start at zero?
  \end{alertblock}
\end{frame}



% SLIDE 19
\begin{frame}{Anatomy of a Bar Graph}
  \centering
  \lectureimage[0.8\linewidth]{fig_bar_anatomy.png}
\end{frame}

% SLIDE 20
\begin{frame}{Pie Charts: Parts of One Whole}
  A \textbf{pie chart} divides a single whole into categorical slices.
  \vspace{0.1cm}
  Use a pie chart only when:
  \begin{enumerate}
    \item all categories belong to the \textbf{same group};
    \item categories are \textbf{mutually exclusive};
    \item categories represent the \textbf{entire sample population}; and
    \item there are only a few categories.
  \end{enumerate}
  \begin{block}{Practical Rule}
    A bar graph is usually the safer default because comparing lengths is easier than comparing angles and areas.
  \end{block}
\end{frame}



% SLIDE 21
\begin{frame}{Anatomy of a Pie Chart}
  \centering
  \lectureimage[0.82\linewidth]{fig_pie_anatomy.png}
\end{frame}

% SLIDE 22
\begin{frame}{Bar Graph vs. Pie Chart}
  \begin{columns}[c]
    \begin{column}{0.58\textwidth}
      \centering
      \lectureimage{fig_bar_vs_pie.png}
    \end{column}
    \begin{column}{0.39\textwidth}
      Both displays show the same blood-type percentages.
      \begin{itemize}
        \item The \textbf{pie chart} emphasizes part-to-whole structure.
        \item The \textbf{bar graph} makes small differences easier to compare.
      \end{itemize}
      \begin{block}{Recommendation}
        Use a pie chart only when the part-to-whole relationship is the main message.
      \end{block}
    \end{column}
  \end{columns}
\end{frame}

% =============================================================
% PART IV: MISLEADING OR INAPPROPRIATE GRAPHS
% =============================================================
\section{Misleading Graphs}

% SLIDE 23
\begin{frame}{What Can Go Wrong?}
  A graph can look polished and still communicate the data poorly.
  \vspace{0.2cm}
  In this chapter, we will watch for three common problems:
  \begin{enumerate}
    \item using a \textbf{line graph for nominal categories};
    \item \textbf{truncating the vertical axis} of a bar graph; and
    \item making a \textbf{pie chart from percentages with different denominators}.
  \end{enumerate}
  \vspace{0.2cm}
  \begin{alertblock}{Statistical Literacy}
    Creating a graph is only half the job. You must also be able to judge whether a graph explains \textbf{what data can't, is appropriate and without bias}.
  \end{alertblock}
\end{frame}

% SLIDE 24
\begin{frame}{Using a Line Graph for Nominal Categories}
  \begin{columns}[c]
    \begin{column}{0.58\textwidth}
      \centering
      \lectureimage{fig_wrong_line_categorical.png}
    \end{column}
    \begin{column}{0.39\textwidth}
      Blood type is a \textbf{nominal categorical variable}.
      \begin{itemize}
        \item There is no continuous path from O to A to B to AB.
        \item Connecting categories with lines suggests a relationship that does not exist.
      \end{itemize}
      \begin{block}{Rule}
        Use lines when \textbf{order or continuity matters}, especially time.
      \end{block}
    \end{column}
  \end{columns}
\end{frame}

% SLIDE 25
\begin{frame}{Truncating the Axis of a Bar Graph}
  \begin{columns}[c]
    \begin{column}{0.58\textwidth}
      \centering
      \lectureimage{fig_truncated_bar_axis.png}
    \end{column}
    \begin{column}{0.39\textwidth}
      Catalyst A: \textbf{78\% yield}\\
      Catalyst B: \textbf{82\% yield}
      \vspace{0.1cm}
      \begin{itemize}
        \item Actual difference: \textbf{4 percentage points}.
        \item Starting the axis near 75\% makes the bars appear dramatically different.
      \end{itemize}
      \begin{alertblock}{Bar-Graph Rule}
        Because bar length represents magnitude, the vertical axis should \textbf{usually begin at zero}.
      \end{alertblock}
    \end{column}
  \end{columns}
\end{frame}

% SLIDE 26
\begin{frame}{A Pie Chart with Different Denominators}
  \begin{columns}[c]
    \begin{column}{0.58\textwidth}
      \centering
      \lectureimage{fig_wrong_pie_denominators.png}
    \end{column}
    \begin{column}{0.39\textwidth}
      Suppose the percentages sleeping at least 7 hours are:
      \begin{itemize}
        \item Year 1: 72\%
        \item Year 2: 65\%
        \item Year 3: 61\%
        \item Year 4: 58\%
      \end{itemize}
      Each percentage comes from a \textbf{different group}.
      \begin{block}{Correct Choice}
        Compare the group rates with a \textbf{bar graph}, not a pie chart.
      \end{block}
    \end{column}
  \end{columns}
\end{frame}

% SLIDE 27
\begin{frame}{A Quick Fair-Graph Checklist}
  Before trusting a graph, ask:
  \begin{itemize}
    \item Does the graph type match the \textbf{variable type}?
    \item Are the \textbf{axes} clearly labeled?
    \item Are the \textbf{units} shown?
    \item Does a bar graph begin at \textbf{zero}?
    \item Are percentages using the \textbf{same denominator}?
    \item Is anything important hidden by the \textbf{scale, bin width, or ordering}?
    \item Does the title describe the data without exaggerating the conclusion?
  \end{itemize}
  \begin{alertblock}{Remember}
    A graph is an argument about the data. Design choices can make that argument more honest---or more misleading.
  \end{alertblock}
\end{frame}

% =============================================================
% PART V: HISTOGRAMS AND SOCS
% =============================================================
\section{Histograms and SOCS}

% SLIDE 26
\begin{frame}{From Numerical Measurements to a Distribution}
  Quantitative datasets may contain many distinct values:
  \[
    17.8,\;18.1,\;18.3,\;18.4,\;18.7,\ldots
  \]
  A separate bar for every exact value would quickly become difficult to read.
  \vspace{0.2cm}
  \begin{block}{The Histogram Solution}
    Group nearby numerical values into intervals called \textbf{bins} or \textbf{classes} and count how many observations fall in each bin.
  \end{block}
\end{frame}

% SLIDE 29
\begin{frame}{How a Histogram Works}
  A histogram:
  \begin{enumerate}
    \item divides the numerical scale into intervals;
    \item counts how many observations fall in each interval; and
    \item draws bars whose heights show those frequencies.
  \end{enumerate}
  \vspace{0.2cm}
  \begin{alertblock}{Crucial Difference from a Bar Graph}
    Histogram bars are adjacent to eachother because the horizontal axis is a continuous numerical scale.
  \end{alertblock}
\end{frame}



% SLIDE 30
\begin{frame}{Anatomy of a Histogram}
  \centering
  \lectureimage[0.92\linewidth]{fig_hist_anatomy.png}
\end{frame}

% SLIDE 31
\begin{frame}{Bar Graph vs. Histogram}
  \begin{table}
    \centering
    \begin{tabular}{lll}
      \toprule
      \textbf{Feature} & \textbf{Bar Graph} & \textbf{Histogram} \\
      \midrule
      Variable type & Categorical & Quantitative \\
      Horizontal axis & Category labels & Numerical intervals \\
      Bars & Separated & Adjacent \\
      Order & May be rearranged & Must follow numerical order \\
      Main purpose & Compare categories & Show distribution shape \\
      \bottomrule
    \end{tabular}
  \end{table}
  \vspace{0.2cm}
  \centering
  \textbf{Do not choose between them based only on how the graph looks. Start with the variable type.}
\end{frame}

% SLIDE 29
\begin{frame}{Example 1.2: Body Lengths of 44 Great White Sharks}
  Marine biologists measured the body lengths of $n=44$ great white sharks.
  \vspace{0.15cm}
  \begin{block}{A Few of the Measurements (feet)}
    \centering
    $18.7,\;12.3,\;18.6,\;16.4,\;15.7,\;18.3,\;14.6,\;15.8,\;14.9,\;17.6,\ldots$
  \end{block}
  \vspace{0.15cm}
  Rather than draw 44 separate bars, we group nearby values into length classes.
  \begin{itemize}
    \item Example class: 13.1 to 15.0 ft
    \item Number of sharks in that class: 12
    \item Approximate percentage: 27\%
  \end{itemize}
\end{frame}

% SLIDE 30
\begin{frame}{Example 1.2: Shark Length Frequency Table}
  \begin{table}
    \centering
    \small
    \begin{tabular}{lrr}
      \toprule
      \textbf{Length Class (feet)} & \textbf{Count} & \textbf{Approx. \%} \\
      \midrule
      9.1 to 11.0 & 1 & 2\% \\
      11.1 to 13.0 & 5 & 11\% \\
      13.1 to 15.0 & 12 & 27\% \\
      15.1 to 17.0 & 15 & 34\% \\
      17.1 to 19.0 & 8 & 18\% \\
      19.1 to 21.0 & 2 & 5\% \\
      21.1 to 23.0 & 1 & 2\% \\
      \midrule
      \textbf{Total} & \textbf{44} & \textbf{99\%*} \\
      \bottomrule
    \end{tabular}
  \end{table}
  \small\centering
  *Percentages total 99\% because of rounding.
\end{frame}

% SLIDE 31
\begin{frame}{Example 1.2: The Same Shark Data, Different Bins}
  \begin{columns}[c]
    \begin{column}{0.59\textwidth}
      \centering
      \lectureimage{fig_sharks_hist.png}
    \end{column}
    \begin{column}{0.38\textwidth}
      Both histograms use the same 44 shark lengths.
      \begin{itemize}
        \item The main cluster remains around the middle lengths.
        \item Changing bin width changes the amount of visible detail.
        \item The smallest and largest observations become more or less obvious.
      \end{itemize}
    \end{column}
  \end{columns}
\end{frame}

% SLIDE 32
\begin{frame}{The Bin Width Problem}
  \begin{columns}[c]
    \begin{column}{0.60\textwidth}
      \centering
      \lectureimage{fig_hist_binwidth.png}
    \end{column}
    \begin{column}{0.37\textwidth}
      \begin{itemize}
        \item \textbf{Too few bins:} important features are smoothed away.
        \item \textbf{Too many bins:} random noise can look meaningful.
        \item \textbf{Reasonable width:} shows structure without manufacturing a story.
      \end{itemize}
      \begin{block}{Good Practice}
        Inspect a few reasonable choices and check whether the main interpretation stays similar.
      \end{block}
    \end{column}
  \end{columns}
\end{frame}

% SLIDE 33
\begin{frame}{How to Describe a Quantitative Distribution: SOCS}
  When asked to describe a histogram or dotplot, use \textbf{SOCS}:
  \vspace{0.25cm}
  \begin{description}[Outliers:]
    \item[\textbf{S -- Shape:}] Is it symmetric, skewed left, skewed right, unimodal, or bimodal?
    \item[\textbf{O -- Outliers:}] Are any observations far from the overall pattern?
    \item[\textbf{C -- Center:}] Where is a typical or middle value?
    \item[\textbf{S -- Spread:}] How variable are the observations? What is the approximate range?
  \end{description}
\end{frame}

% SLIDE 34
\begin{frame}{Shape: Symmetry, Skewness and Peaks}
  \begin{columns}[c]
    \begin{column}{0.57\textwidth}
      \centering
      \lectureimage{fig_hist_socs.png}
    \end{column}
    \begin{column}{0.40\textwidth}
      Common shapes include:
      \begin{itemize}
        \item \textbf{symmetric};
        \item \textbf{skewed right};
        \item \textbf{skewed left};
        \item \textbf{unimodal};
        \item \textbf{bimodal}.
      \end{itemize}
      \begin{block}{Important}
        Skewness is named for the direction of the \textbf{long tail}, not the location of the tallest bar.
      \end{block}
    \end{column}
  \end{columns}
\end{frame}

% SLIDE 35
\begin{frame}{Outliers, Center and Spread}
  \begin{columns}[T]
    \begin{column}{0.31\textwidth}
      \begin{block}{Outliers}
        Values far from the overall pattern. They may be genuine or may indicate an error.
      \end{block}
    \end{column}
    \begin{column}{0.31\textwidth}
      \begin{block}{Center}
        A typical or middle value. For now, estimate visually; Chapter 2 introduces mean and median.
      \end{block}
    \end{column}
    \begin{column}{0.31\textwidth}
      \begin{block}{Spread}
        How much the values vary. A narrow distribution has less spread than a wide one.
      \end{block}
    \end{column}
  \end{columns}
  \vspace{0.3cm}
  \begin{alertblock}{Never Delete an Observation Just Because It Looks Unusual}
    Investigate the value, its units, and how it was recorded before deciding what to do with it.
  \end{alertblock}
\end{frame}

% SLIDE 36
\begin{frame}{Example 1.3: SOCS for Shark Lengths}
  \begin{columns}[c]
    \begin{column}{0.55\textwidth}
      \centering
      \lectureimage{fig_sharks_hist.png}
    \end{column}
    \begin{column}{0.42\textwidth}
      \begin{itemize}
        \item \textbf{Shape:} roughly unimodal and fairly symmetric through the center.
        \item \textbf{Outliers:} 9.4 ft and 22.8 ft appear separated from much of the data.
        \item \textbf{Center:} approximately 15--16 ft.
        \item \textbf{Spread:} 9.4 ft to 22.8 ft.
      \end{itemize}
      \begin{block}{Goal}
        At this stage, learn to \textbf{read the distribution}; exact numerical summaries come next chapter.
      \end{block}
    \end{column}
  \end{columns}
\end{frame}

% =============================================================
% PART VI: DOTPLOTS
% =============================================================
\section{Dotplots}

% SLIDE 37
\begin{frame}{Dotplots: Showing Individual Values}
  A histogram groups observations into bins. For a smaller dataset, we often want to see every value.
  \vspace{0.15cm}
  A \textbf{dotplot} places one dot for each observation above a numerical scale.
  \begin{itemize}
    \item Individual values remain visible.
    \item Repeated values can be stacked.
    \item Clusters, gaps, center, spread, and outliers are easy to inspect.
    \item Small groups can be compared directly.
  \end{itemize}
  \begin{block}{When?}
    The sample is small or moderate and the individual observations are still readable.
  \end{block}
\end{frame}


% SLIDE 39
\begin{frame}{Anatomy of a Dotplot}
  \centering
  \lectureimage[0.90\linewidth]{fig_dotplot_anatomy.png}
\end{frame}

% SLIDE 40
\begin{frame}{Example 1.4: Plant Height With and Without Fertilizer}
  \begin{columns}[c]
    \begin{column}{0.59\textwidth}
      \centering
      \lectureimage{fig_plant_dotplot.png}
    \end{column}
    \begin{column}{0.38\textwidth}
      Each dot represents one plant.
      \begin{itemize}
        \item Fertilizer plants tend to be taller.
        \item The groups still overlap.
        \item No single point is extremely separated from the rest.
      \end{itemize}
      \begin{alertblock}{Careful Wording}
        The graph describes the \textbf{observed samples}. By itself, it does not prove what will happen for every future plant.
      \end{alertblock}
    \end{column}
  \end{columns}
\end{frame}

% SLIDE 41
\begin{frame}{Dotplot vs. Histogram}
  \begin{table}
    \centering
    \begin{tabular}{ll}
      \toprule
      \textbf{Dotplot} & \textbf{Histogram} \\
      \midrule
      Shows individual observations & Groups observations into bins \\
      Excellent for smaller samples & Excellent for larger samples \\
      Repeated values remain visible & Overall shape is easier with dense data \\
      Sample size can often be counted & Exact values are not visible \\
      \bottomrule
    \end{tabular}
  \end{table}
  \vspace{0.25cm}
  \begin{block}{No Magic Cutoff}
    There is no universal rule that ``30 observations means histogram.'' Choose the display that communicates the data most clearly.
  \end{block}
\end{frame}

% =============================================================
% PART VII: TIME PLOTS
% =============================================================
\section{Time Plots}

% SLIDE 42
\begin{frame}{Time Plots: When Order Matters}
  A histogram or dotplot tells us \textbf{what values occurred}. It does not tell us \textbf{when they occurred}.
  \vspace{0.15cm}
  When chronological order matters (such as time series), use a \textbf{time plot}:
  \begin{itemize}
    \item \textbf{$x$-axis:} time;
    \item \textbf{$y$-axis:} measured variable;
    \item points are commonly connected in chronological order.
  \end{itemize}
  \vspace{0.15cm}
  Look for:
  \begin{itemize}
    \item \textbf{trend}
    \item \textbf{cycle}
    \item \textbf{sudden change}
    \item \textbf{unusual observation}
  \end{itemize}
\end{frame}


% SLIDE 43
\begin{frame}{Anatomy of a Time Plot}
  \centering
  \lectureimage[0.90\linewidth]{fig_timeplot_anatomy.png}
\end{frame}

% SLIDE 44
\begin{frame}{Example 1.5: A Repeating Daily Pattern}
  \begin{columns}[c]
    \begin{column}{0.59\textwidth}
      \centering
      \lectureimage{fig_timeplot.png}
    \end{column}
    \begin{column}{0.38\textwidth}
      The measurements rise and fall in a repeating pattern over approximately 24 hours.
      \begin{itemize}
        \item A histogram could summarize the temperatures.
        \item But it would hide \textbf{when} the high and low values occurred.
      \end{itemize}
      \begin{block}{Time-Plot Advantage}
        Time order makes cycles and trends visible.
      \end{block}
    \end{column}
  \end{columns}
\end{frame}

% SLIDE 45
\begin{frame}{Example 1.6: A Histogram Can Hide Instrument Drift}
  \begin{columns}[c]
    \begin{column}{0.60\textwidth}
      \centering
      \lectureimage{fig_time_order_drift.png}
    \end{column}
    \begin{column}{0.37\textwidth}
      The two panels contain the \textbf{same measurements}.
      \begin{itemize}
        \item Histogram: what concentration values occurred?
        \item Time plot: when did they occur?
        \item The time plot reveals a gradual upward drift.
      \end{itemize}
      \begin{alertblock}{Golden Rule}
        When time order is meaningful, \textbf{plot against time before ignoring the order}.
      \end{alertblock}
    \end{column}
  \end{columns}
\end{frame}

% =============================================================
% PART VIII: OUTLIERS AND DATA QUALITY
% =============================================================
\section{Data Quality}

% SLIDE 43
\begin{frame}{Outliers and Data Quality: Always Investigate}
  When you see an unusual value, ask:
  \begin{enumerate}
    \item Was the value entered correctly?
    \item Are the units correct?
    \item Was the instrument working correctly?
    \item Is the value scientifically plausible?
    \item Could it be a genuine but unusual observation?
  \end{enumerate}
  \vspace{0.2cm}
  \begin{alertblock}{Important}
    An outlier is \textbf{not automatically an error}. Do not remove data simply because it looks unusual.
  \end{alertblock}
\end{frame}

% SLIDE 44
\begin{frame}{Example 1.7: A Data-Entry Error}
  \begin{columns}[c]
    \begin{column}{0.59\textwidth}
      \centering
      \lectureimage{fig_data_entry.png}
    \end{column}
    \begin{column}{0.38\textwidth}
      Most body temperatures are near 37$^\circ$C, but one value is recorded as \textbf{370$^\circ$C}.
      \begin{itemize}
        \item The graph makes the error obvious.
        \item Without plotting first, the typo could badly distort numerical summaries.
      \end{itemize}
    \end{column}
  \end{columns}
\end{frame}

% SLIDE 45
\begin{frame}{Common Data Problems That Graphs Can Reveal}
  \begin{table}
    \centering
    \small
    \begin{tabularx}{\textwidth}{>{\bfseries}l X X}
      \toprule
      \textbf{Problem} & \textbf{Example} & \textbf{Possible Visual Signature} \\
      \midrule
      Typing error & 19.4 entered as 194 & One extreme observation \\
      Wrong unit & Some masses in mg instead of g & Two separated groups or a huge gap \\
      Missing-value code & Missing recorded as $-99$ & Strange stack at $-99$ \\
      Instrument drift & Calibration changes over days & Gradual trend in a time plot \\
      Duplicated record & Same row copied twice & Unusually repeated points \\
      \bottomrule
    \end{tabularx}
  \end{table}
  \vspace{0.15cm}
  \centering
  \textbf{Graphs are tools for checking data quality, not only for presenting results.}
\end{frame}

% =============================================================
% PART IX: PRACTICE AND SYNTHESIS
% =============================================================
\section{Practice and Summary}

% SLIDE 46
\begin{frame}{Apply Your Knowledge 1.1: Breakfast Cereals}
  Consider a dataset describing 77 breakfast cereals:
  \begin{table}
    \centering
    \small
    \begin{tabular}{lcccc}
      \toprule
      \textbf{Brand} & \textbf{Manufacturer} & \textbf{Type} & \textbf{Calories} & \textbf{Sugar (g)} \\
      \midrule
      All Bran & K & C & 70 & 5 \\
      All Bran Fiber & K & C & 50 & 0 \\
      Almond Delight & R & C & 110 & 8 \\
      Apple Cheerios & G & C & 110 & 10 \\
      Apple Jacks & K & C & 110 & 14 \\
      \bottomrule
    \end{tabular}
  \end{table}
  \small
  \textbf{Questions:}
  \begin{enumerate}
    \item What are the individuals?
    \item Which variables are categorical? Which are quantitative?
    \item Which graph would you use for \texttt{Manufacturer}?
    \item Which graph would you use for \texttt{Sugar} across all 77 cereals?
  \end{enumerate}
\end{frame}

% SLIDE 47
\begin{frame}{Apply Your Knowledge 1.1: Answers}
  \begin{itemize}
    \item \textbf{Individuals:} the 77 cereal brands.
    \item \textbf{Categorical:} Manufacturer and Type.
    \item \textbf{Quantitative:} Calories and Sugar.
    \item \textbf{Manufacturer:} bar graph.
    \item \textbf{Sugar for 77 cereals:} histogram is a good starting choice.
  \end{itemize}
  \vspace{0.2cm}
  \begin{block}{Reasoning Pattern}
    Do not memorize graph names in isolation. Start by identifying the \textbf{individual}, then the \textbf{variable type}, then the \textbf{question}.
  \end{block}
\end{frame}

% SLIDE 48
\begin{frame}{Apply Your Knowledge 1.2: Choose the Graph}
  \begin{table}
    \centering
    \small
    \begin{tabularx}{\textwidth}{X X l}
      \toprule
      \textbf{Question} & \textbf{Structure} & \textbf{Best Starting Graph} \\
      \midrule
      Which blood type is most common? & Categorical & Bar graph \\
      How are four categories divided within one whole? & Categorical, part-to-whole & Pie or bar graph \\
      What is the distribution of 150 plant heights? & Quantitative & Histogram \\
      What are the exact reaction times of 12 participants? & Quantitative, small sample & Dotplot \\
      Is river temperature changing throughout the day? & Quantitative over time & Time plot \\
      \bottomrule
    \end{tabularx}
  \end{table}
\end{frame}

% SLIDE 49
%\begin{frame}{A Five-Step Approach for Exam Questions}
%  When given a new dataset or graph-selection problem:
%  \begin{enumerate}
%    \item \textbf{Identify the individual:} What does one row represent?
%    \item \textbf{Identify the variable:} What characteristic was recorded?
%    \item \textbf{Classify the variable:} Categorical or quantitative? If quantitative, discrete or continuous?
%    \item \textbf{Ask whether time order matters.}
%    \item \textbf{Choose the graph} that answers the question most clearly.
%  \end{enumerate}
%  \vspace{0.2cm}
%  \begin{alertblock}{Then Interpret}
%    Do not stop after naming the graph. State the main pattern that the graph is meant to reveal.
%  \end{alertblock}
%\end{frame}

% SLIDE 50
\begin{frame}{Master Guide: Which Graph Should You Use?}
  \begin{table}
    \centering
    \scriptsize
    \begin{tabularx}{\textwidth}{X l X X}
      \toprule
      \textbf{Data Situation} & \textbf{Graph} & \textbf{What to Describe} & \textbf{Common Mistake} \\
      \midrule
      Categorical counts / percentages & Bar & Largest/smallest categories; comparisons & Connecting categories with a line \\
      Parts of one whole & Pie & Dominant slices; part-to-whole & Using separate group rates as slices \\
      Small quantitative sample & Dotplot & Individual values; SOCS & Hiding a tiny sample in bins \\
      Larger quantitative sample & Histogram & SOCS; overall shape & Poor bin choice \\
      Data recorded over time & Time plot & Trend; cycles; sudden changes & Ignoring time order \\
      \bottomrule
    \end{tabularx}
  \end{table}
\end{frame}

% SLIDE 51
%\begin{frame}{What Should You Say When Describing a Graph?}
%  \begin{columns}[T]
%    \begin{column}{0.31\textwidth}
%      \begin{block}{Bar Graph}
%        \begin{itemize}
%          \item largest / smallest categories;
%          \item important comparisons;
%          \item counts or percentages;
%          \item denominator when relevant.
%        \end{itemize}
%      \end{block}
%    \end{column}
%    \begin{column}{0.31\textwidth}
%      \begin{block}{Histogram / Dotplot}
%        Use \textbf{SOCS}:
%        \begin{itemize}
%          \item Shape
%          \item Outliers
%          \item Center
%          \item Spread
%        \end{itemize}
%      \end{block}
%    \end{column}
%    \begin{column}{0.31\textwidth}
%      \begin{block}{Time Plot}
%        Describe:
%        \begin{itemize}
%          \item trend;
%          \item cycles;
%          \item sudden changes;
%          \item unusual observations.
%        \end{itemize}
%      \end{block}
%    \end{column}
%  \end{columns}
%\end{frame}

% SLIDE 52
%\begin{frame}{Chapter 1 Summary: The Core Rules}
%  \begin{enumerate}
%    \item Know what \textbf{one row represents}.
%    \item Classify the variable before choosing a graph.
%    \item Use \textbf{bar graphs} for categorical comparisons; pie charts only for simple parts of one whole.
%    \item Use \textbf{histograms} for the shape of larger quantitative datasets and inspect the bin width.
%    \item Use \textbf{dotplots} when individual values remain readable.
%    \item Use \textbf{time plots} when chronological order matters.
%    \item Describe quantitative distributions using \textbf{SOCS}.
%    \item Check axes, scales, denominators, labels, and other design choices.
%    \item \textbf{Plot before calculating} so errors and unusual observations are visible.
%  \end{enumerate}
%\end{frame}

% SLIDE 53
%\begin{frame}{Final Graph Checklist}
%  Before interpreting or presenting a graph, ask:
%  \begin{itemize}
%    \item What are the \textbf{individuals}?
%    \item What \textbf{variable} is being displayed?
%    \item Is it \textbf{categorical or quantitative}?
%    \item Does \textbf{time order} matter?
%    \item Is this the \textbf{right graph} for the variable and question?
%    \item Are the \textbf{title, axes, and units} clear?
%    \item Could the \textbf{scale or binning} exaggerate or hide a pattern?
%    \item What is the \textbf{main pattern} visible in the data?
%  \end{itemize}
%\end{frame}



% =============================================================
% PRACTICE: FIND THE ISSUE WITH THE GRAPH
% =============================================================
\section{Practice: Spot the Problem}

% SLIDE 54 -- EASY
\begin{frame}{Practice 1 (Easy): What Is Wrong with This Graph?}
  \begin{columns}[c]
    \begin{column}{0.57\textwidth}
      \centering
      \lectureimage[0.98\linewidth]{fig_practice_easy_wrong_line.png}
    \end{column}
    \begin{column}{0.40\textwidth}
      \begin{block}{Your Task}
        The graph shows students' \textbf{favourite study drink}. Identify the main problem with the display.
      \end{block}
      \textbf{Ask yourself:}
      \begin{itemize}
        \item What type of variable is \emph{drink}?
        \item Does connecting the categories imply something that is not real?
      \end{itemize}
      \pause
      \begin{alertblock}{Answer}
        \small \textbf{Drink is nominal categorical.} Connecting the categories suggests an order or continuous movement that does not exist. Use a \textbf{bar graph}.
      \end{alertblock}
    \end{column}
  \end{columns}
\end{frame}

% SLIDE 55 -- MEDIUM
\begin{frame}{Practice 2 (Medium): Is the Difference Really That Large?}
  \begin{columns}[c]
    \begin{column}{0.57\textwidth}
      \centering
      \lectureimage[0.98\linewidth]{fig_practice_medium_truncated_bar.png}
    \end{column}
    \begin{column}{0.40\textwidth}
      \begin{block}{Your Task}
        Treatment A has a success rate of \textbf{78\%}; Treatment B has a success rate of \textbf{82\%}. What makes this graph misleading?
      \end{block}
      \textbf{Look carefully at:}
      \begin{itemize}
        \item the vertical axis
        \item the visual size of the difference
        \item the actual numerical difference
      \end{itemize}
      \pause
      \begin{alertblock}{Answer}
        \small The axis begins near \textbf{75\%}, not 0, which exaggerates the \textbf{4-point difference}. Bar charts should generally start at zero.
      \end{alertblock}
    \end{column}
  \end{columns}
\end{frame}

% SLIDE 56 -- HARD
\begin{frame}{Practice 3 (Hard): What Important Information Is Missing?}
  \begin{columns}[c]
    \begin{column}{0.57\textwidth}
      \centering
      \lectureimage[0.98\linewidth]{fig_practice_hard_histogram_time.png}
    \end{column}
    \begin{column}{0.40\textwidth}
      \begin{block}{Research Question}
        A temperature was recorded every hour for 24 hours. The researcher wants to know:\\[0.08cm]
        \emph{Did the reaction vessel temperature drift during the experiment?}
      \end{block}
      \textbf{What important information has the graph discarded?}
      \pause
      \begin{alertblock}{Answer}
        \small The histogram shows the temperature distribution but removes \textbf{time order}. It cannot reveal drift, cycles, or sudden changes. Plot \textbf{temperature versus time} first.
      \end{alertblock}
    \end{column}
  \end{columns}
\end{frame}

% SLIDE 57
\begin{frame}{Chapter 1 Summary: The Core Rules}
  \begin{enumerate}
    \item Know what \textbf{one row represents}.
    \item Classify the variable before choosing a graph.
    \item Use \textbf{bar graphs} for categorical comparisons; pie charts only for simple parts of one whole.
    \item Use \textbf{histograms} for the shape of larger quantitative datasets and inspect the bin width.
    \item Use \textbf{dotplots} when individual values remain readable.
    \item Use \textbf{time plots} when chronological order matters.
    \item Describe quantitative distributions using \textbf{SOCS}.
    \item Check axes, scales, denominators, labels, and other design choices.
    \item \textbf{Plot before calculating} so errors and unusual observations are visible.
  \end{enumerate}
\end{frame}

% SLIDE 58
\begin{frame}{Looking Ahead: Chapter 2}
  In Chapter 1, we learned to \textbf{see} center and spread.
  \vspace{0.2cm}
  In Chapter 2, we will learn to \textbf{calculate} them.
  \[
    \text{Mean} \qquad \text{Median}
  \]
  and measures of variability such as
  \[
    \text{Range} \qquad \text{Interquartile Range} \qquad \text{Standard Deviation}.
  \]
  \vspace{0.25cm}
  \begin{block}{The Same Principle Continues}
    \textbf{Numbers and graphs should support each other.}
  \end{block}
\end{frame}

\end{document}
